% !PRINTED 378
\weberpage{378}
Die Zahlen $a$, $b$, $c$ sind durch (4) nur bis auf ein gemeinsames Vorzeichen bestimmt.

Die Discriminante von (3) ist
\weq{6}{D = b^2 + 4ac = 4c^2 y^2 d,}
und soll auch die \gesperrt{Discriminante der Irrationalzahl} $\omega$ heissen. $D$ ist also eine ganze Zahl, die im Vorzeichen mit $d$ übereinstimmt und die durch $d$ theilbar ist; denn ist $p$ ein Primtheiler von $d$, der also in $d$ nach Voraussetzung nur einmal aufgeht, so kann $p$ nicht im Nenner von $2cy$ aufgehen und muss also nach (6) in $D$ aufgehen. Also ist auch $2cy$ eine ganze Zahl und der Quotient $D:d$ ist eine Quadratzahl.

Nach (4) und (6) lässt sich $\omega$ und $\omega'$ so darstellen:
\weq{7}{\omega = \frac{\sqrt{D}+b}{2c}, \quad \omega' = \frac{-\sqrt{D}+b}{2c},}
oder was dasselbe ist:
\weq{8}{\omega = \frac{2a}{\sqrt{D}-b}, \quad \omega' = \frac{-2a}{\sqrt{D}+b},}
wenn das Vorzeichen der Wurzel aus
\weq{9}{\sqrt{D} = 2cy\sqrt{d}}
bestimmt wird. Da eine gerade Quadratzahl durch 4 theilbar ist, eine ungerade durch 4 getheilt, den Rest 1 lässt, so ergiebt sich aus (6), dass
\weq{10}{D \equiv 0 \quad \text{oder} \quad \equiv 1 \ (\mathrm{mod}\ 4)}
sein muss.

Wir wollen noch untersuchen, wie sich die Zahlen $a$, $b$, $c$, $D$ ändern, wenn wir von $\omega$ zu einer äquivalenten Zahl $\omega_1$ übergehen. Es sei also
\weq{11}{\omega_1 = \frac{\alpha\omega+\beta}{\gamma\omega+\delta}, \quad \omega = \frac{\delta\omega_1-\beta}{-\gamma\omega_1+\alpha_1} \quad \cdot \quad \alpha\delta - \beta\gamma = \varepsilon.} % [sic] denominator prints $\alpha_1$, not $\alpha$; source uses a centred-dot separator, not a comma, before the third identity
Aus (5) ergiebt sich für $\omega_1$ die quadratische Gleichung
\weq{12}{c_1\omega_1^2 = a_1 + b_1\omega_1,}
wenn
\weq{13}{\begin{aligned}
-a_1 &= -a\alpha^2 + b\alpha\beta + c\beta^2\\
b_1 &= -2a\alpha\gamma + b(\alpha\delta+\beta\gamma) + 2c\beta\delta,\\
c_1 &= -a\gamma^2 + b\gamma\delta + c\delta^2
\end{aligned}}
gesetzt wird. Aus (13) erhält man durch Auflösung nach $a$, $b$, $c$:
%CONT>>
